\documentclass[11pt]{article}
\usepackage[margin=1in]{geometry}
\usepackage[utf8]{inputenc}
\usepackage[T1]{fontenc}
\usepackage{lmodern}
\usepackage{amsmath,amssymb,amsthm}
\usepackage[title]{appendix}
\usepackage[colorlinks=true,linkcolor=blue,citecolor=blue,urlcolor=blue,pdfauthor={P. Graham},pdftitle={A counterexample to the very sparse greedy conjecture for random hitting sets}]{hyperref}
\DeclareUnicodeCharacter{2265}{\ensuremath{\geq}}
\DeclareUnicodeCharacter{2264}{\ensuremath{\leq}}
\DeclareUnicodeCharacter{2227}{\ensuremath{\wedge}}
\DeclareUnicodeCharacter{2200}{\ensuremath{\forall}}
\DeclareUnicodeCharacter{211D}{\ensuremath{\mathbb{R}}}
\DeclareUnicodeCharacter{2115}{\ensuremath{\mathbb{N}}}
\newtheorem{theorem}{Theorem}
\newtheorem{corollary}[theorem]{Corollary}
\newtheorem{lemma}[theorem]{Lemma}
\theoremstyle{remark}
\newtheorem*{remark}{Remark}
\newcommand{\Gu}{G_{\mathrm u}}
\setlength{\emergencystretch}{3em}
\title{A counterexample to the very sparse greedy conjecture for random hitting sets}
\author{P. Graham\\ \small\url{https://orcid.org/0009-0007-8493-1505}}
\date{November 1, 2026}
\begin{document}
\maketitle
\begin{abstract}
Let $A$ be a random $m\times n$ zero-one matrix with independent
Bernoulli$(p)$ entries, whose rows are to be hit by a set of columns.
Arpino, Dmitriev and Grometto proved that with high probability the
linear programming (LP) relaxation, the fractional hitting set, is within
a constant factor of the optimum when $mp\lesssim\log n$, and the greedy
algorithm is within a constant factor of the optimum when $p\ge n^{-\delta}$
with $\delta<1/2$. They conjectured that in the very sparse regime
$mp\to0$ the ratio of the greedy value to the LP value tends to one. We
show that this fails. For $m=k$, $n=k^{10}$ and $p=1/(10k^4)$, with the
explicit constant $c=1/(196608\cdot10^9)$, the greedy algorithm that
breaks ties by the leftmost column, as in their analysis, uses at least
$(1+c/8)$ times the LP value with probability at least $c/4$ for every
$k\ge10^{20}$; these parameters satisfy their Assumption~2. The same holds
with probability at least $c/8$ when ties are broken uniformly at random.
So the ratio does not converge to one in probability. The obstruction is
a path of three triples whose middle triple greedy takes first; exchanging
it for its two ends saves half a column per path. The proofs are elementary
counting arguments, checked in Lean from the definitions.

This work has not yet been reviewed by experts.
The proofs, the Lean formalization and the text were produced by AI models
(from OpenAI and Anthropic), working under the author's direction.
\end{abstract}
\medskip
\noindent\textit{Keywords:} hitting set, set cover, greedy algorithm,
linear programming relaxation, random hypergraph, formal proof.

\noindent\textit{MSC 2020:} 05C80, 68W40, 90C05, 60C05.

\section{Introduction}

\subsection{Main results}
Let $A$ be an $m\times n$ matrix with entries in $\{0,1\}$. Its rows are
the sets to be hit and its columns the elements; write
$A_j=\{i:A_{ij}=1\}$ for the set of rows that column $j$ hits. The matrix
is \emph{feasible} if every row is hit by some column. The greedy
algorithm keeps a set $U$ of uncovered rows, initially all of them. While
$U\ne\emptyset$, it chooses a column $j$ maximizing the \emph{gain}
$|A_j\cap U|$, removes $A_j$ from $U$, and counts the choice. We write
$G(A)$ for the count when ties are broken by the smallest column index.
On a feasible matrix it stops with a set of $G(A)$ columns hitting every
row. The LP value is
\[
 L(A)=\min\Bigl\{\sum_{j=1}^n x_j : x\in[0,1]^n,\ \sum_{j:\,i\in A_j}x_j\ge1
 \text{ for every row } i\Bigr\}.
\]
Always $L(A)\le G(A)$ on feasible matrices.

Arpino, Dmitriev and Grometto \cite{ADG} study a random matrix whose
entries are independent Bernoulli$(p)$ variables, with $m=m(n)$ and
$p=p(n)$. Their Assumption~2 \cite[p.~30:3]{ADG} requires
$c_0n^{c_0}\le m\le C_0n^{C_0}$ for some $c_0,C_0>0$, and $n^{-\delta}\le p\le1/2$ for
some $\delta\in(0,1)$, for all large $n$. Their Conjecture~12
\cite[Section~6.2, p.~30:13]{ADG}, labeled ``Very sparse'', reads: ``For
$mp\ll1$, $\mathrm{val}_{\mathrm{Gr}}/\mathrm{val}_{\mathrm{LP}}\to1$.''
Here $\mathrm{val}_{\mathrm{Gr}}$ is the output of their Algorithm~1
\cite[p.~30:7]{ADG}, the greedy algorithm above, and
$\mathrm{val}_{\mathrm{LP}}$ is the LP value \cite[equation~(1.2)]{ADG}.
In the proof of their Theorem~10 \cite[p.~30:11]{ADG} ties are broken by
``selecting the left-most column in the associated matrix''. The
conjecture first appeared, as Conjecture~1 in Section~7.2, in the preprint
\cite{ADGarxiv} of May 9, 2023, and Dmitriev's thesis restates it as open
\cite[Conjecture~6.1(1)]{Dmitriev}.

Throughout, $c=1/(196608\cdot10^9)$, and for an integer $k\ge1$ we take
\begin{equation}\label{eq:params}
 m=k,\qquad n=k^{10},\qquad p=\frac1{10k^4}.
\end{equation}

\begin{theorem}\label{thm:leftmost}
For every integer $k\ge10^{20}$, with parameters \eqref{eq:params},
\[
 \Pr\Bigl[A\text{ is feasible and }G(A)\ge\bigl(1+\tfrac c8\bigr)L(A)\Bigr]
 \ge\frac c4 .
\]
\end{theorem}

For the second theorem, let $\Gu(A)$ be the number of choices of the greedy
algorithm when, at each step and given everything before it, the column is
chosen uniformly at random among the columns of maximum positive gain.
Columns with equal sets $A_j$ count as different columns.

\begin{theorem}\label{thm:uniform}
For every integer $k\ge10^{20}$, with parameters \eqref{eq:params},
\[
 \Pr\Bigl[A\text{ is feasible and }\Gu(A)\ge\bigl(1+\tfrac c8\bigr)L(A)\Bigr]
 \ge\frac c8 ,
\]
the probability being over the matrix and the choices together.
\end{theorem}

The proof in Section~\ref{sec:uniform} gives the bound $c/4$ of
Theorem~\ref{thm:leftmost}; we state the weaker bound $c/8$, which is the
one in the Lean statement (Section~\ref{sec:lean}).

\begin{corollary}\label{cor:conj}
Let $m(n)=\lfloor n^{1/10}\rfloor$ and $p(n)=n^{-2/5}/10$. This sequence
satisfies Assumption~2 of \cite{ADG}, $m(n)p(n)\to0$, and
$m(k^{10})=k$, $p(k^{10})=1/(10k^4)$. Assign arbitrary real values to the
ratios $G/L$ and $\Gu/L$ on infeasible matrices. Then neither ratio
converges to $1$ in probability as $n\to\infty$.
\end{corollary}

Convergence in probability is implied by almost sure convergence and by
convergence in mean, so the corollary refutes Conjecture~12 under each of
these readings, for the leftmost rule that \cite{ADG} use and for uniform
tie-breaking. The constants are far from optimal; we made no attempt to
improve them. We do not know whether the ratio tends to one for some other
tie-breaking rule, or for the best of all greedy runs, along this
sequence.

\subsection{Prior work}
Lov\'asz \cite{Lovasz} and Johnson \cite{Johnson} analyzed the greedy
algorithm for covering. Lov\'asz showed that ``the ratio of optimal integral
and fractional covers of a hypergraph does not exceed $1+\log d$''
\cite[p.~383]{Lovasz}, with $d$ the maximum degree, as recalled in
\cite[p.~30:2]{ADG}. Slav\'ik bounds the greedy value by
$\ln m-\ln\ln m+0.78$ times the optimum when $m$ elements are to be covered
\cite[Theorem~1]{Slavik} and shows that this is nearly tight
\cite[Theorem~2]{Slavik}. These are worst-case bounds. For the
random model, Corollary~9(ii) of \cite{ADG} states that if $mp\lesssim\log
n$ then $\Pr(\mathrm{IPGAP}\sim1)\ge1-n^{-D}$ for any $D>0$ and $n$ large,
where $\mathrm{IPGAP}=\mathrm{val}_{\mathrm{IP}}/\mathrm{val}_{\mathrm{LP}}$
and $\sim$ means equality up to multiplicative constants; their
Theorem~10 states that under Assumption~2 with $\delta<1/2$,
$\Pr(\mathrm{val}_{\mathrm{Gr}}\sim\mathrm{val}_{\mathrm{IP}})\ge
1-\exp(-\sqrt n)$ for $n$ large. Since $p(n)=n^{-2/5}/10\ge n^{-\delta}$
for large $n$ whenever $\delta>2/5$, these results apply to the sequence of
Corollary~\ref{cor:conj}: greedy is within a constant factor of the LP value
there with high probability. Theorems~\ref{thm:leftmost}
and~\ref{thm:uniform} show that the constant cannot tend to one.

Telelis and Zissimopoulos \cite[Theorem~4, p.~174]{TZ} show, for the same
independent model with fixed $p$, that a simple $O(nm)$ algorithm
approximates random set cover ``within a term of $o(\log m)$ from the
optimum almost everywhere''. Fernandez de la Vega, Paschos and Saad
\cite{FPS} study a greedy algorithm for hitting $m$ random $k$-element
subsets of an $n$-element set, with $k$ fixed and $m,n\to\infty$ in fixed
proportion; their Theorem~1 \cite[p.~131]{FPS} states that its
approximation ratio ``ranges between 1.094 and 1.257'' for $3\le k\le7$
and a range of values of the proportion. See also the later
article of Blot, Fernandez de la Vega, Paschos and Saad \cite{BFPS}.
M\'ezard and Tarzia \cite[Sections~V.C--V.D and Fig.~5]{MT} ran highest-degree
greedy on random hypergraphs in which every variable lies in $L=4$
hyperedges and every hyperedge has $K=6$ variables, with $12288$ variables
and $8192$ hyperedges, and report greedy covers of density
$\rho_{\mathrm{gr}}\simeq0.212$, against $\rho_{\mathrm{cov}}\simeq0.178$
for the minimal cover and $1/K$; on such regular instances the LP value
is $1/K$ times the number of variables.

The isolated bad paths below are paths of three triples, consecutive
triples sharing one row, in a sparse random hypergraph. Greedy matching on sparse random graphs, the setting of
Karp and Sipser \cite{KS} and of Dyer, Frieze and Pittel \cite{DFP}, is
related in spirit, but no matching limit theorem is used here.

We could not read four sources in full: the journal version of
Slav\'ik's paper \cite{Slavik}, whose abstract and references are those of
the conference version cited above; pages 132--138 of \cite{FPS}; Chebolu's
thesis \cite{Chebolu}, which \cite[Section~1]{BB} describes as extending the
analysis of the random greedy matching process \cite{DFP} to sparse random
hypergraphs; and a Zenodo deposit, \emph{Spectral gaps and algorithmic phase
transitions in sparse hypergraphs} (DOI 10.5281/zenodo.17834080), removed
before it could be read, whose surviving abstract concerns coloring,
satisfiability and community detection.

\subsection{Method}
On the event that every column hits at most three rows and every pair of
rows is exactly the set of some column, greedy takes a family of disjoint
triples and then pairs, so its count is determined by the number $K$ of
triples it takes (Lemma~\ref{lem:phases}). The obstruction is an
\emph{isolated bad path}: three triples $\{a,b,c\},\{c,d,e\},\{e,f,g\}$ on
seven rows such that no other column with three rows meets these rows;
the definition in Section~2 adds that each triple is the set of exactly one
column and that the middle column has the smaller index. Leftmost greedy
takes the middle triple, after which the two ends are useless to it; but
the two ends cover six of the seven rows, and the row $d$ joins the pairs
(Lemma~\ref{lem:exchange}). The paths are disjoint, so exchanging every
such middle at once saves half a column per path, up to one rounding.

Counting expected occurrences (a first-moment count) gives order $k$ such
paths in expectation (Lemma~\ref{lem:count}). Their number is at most $m$
on every matrix, so this lower bound on the expectation gives a fixed
fraction of $m$ paths with fixed probability.

For uniform ties, the triples greedy takes have the law of a scan of the
triples in a uniformly random order, which is leftmost greedy on a matrix
with randomly permuted columns; the distribution of $A$ does not change
under this permutation (Section~\ref{sec:uniform}).

The theorems and the corollary are proved in Lean~4 \cite{Lean4} with
Mathlib \cite{Mathlib}, from
definitions of the random matrix, the greedy algorithms and the LP value
(Section~\ref{sec:lean}).

\section{Greedy on the good event}

Let $H$ be the event that $|A_j|\le3$ for every column $j$ and that for
every two rows $a\ne b$ some column has $A_j=\{a,b\}$. If $m\ge2$, then $H$
implies that $A$ is feasible. Gains never increase as $U$ shrinks.

\begin{lemma}\label{lem:phases}
Assume $H$ and $m\ge2$, and break ties by any rule. Greedy first chooses
$K$ columns of gain $3$, whose sets are pairwise disjoint, then columns of
gain $2$ while at least two rows are uncovered, then at most one column
of gain $1$. Hence $G=\lceil(m-K)/2\rceil$, with $G$ the number of
choices.
\end{lemma}

\begin{proof}
Gains are at most $3$ and do not increase. While two rows $a,b$ remain
uncovered, the column with set $\{a,b\}$ has gain $2$, so a choice of gain
$1$ happens only when one row remains. A column of gain $3$ hits three
uncovered rows, so its set is disjoint from those of earlier choices. If
there are $T$ choices of gain $2$ and $S\in\{0,1\}$ of gain $1$, then
$m=3K+2T+S$, so $S\equiv m-3K\pmod2$ and $K+T+S=\lceil(m-K)/2\rceil$.
\end{proof}

An \emph{isolated bad path} consists of seven distinct rows $a,\dots,g$
and three columns with sets $P=\{a,b,c\}$, $M=\{c,d,e\}$, $Q=\{e,f,g\}$
such that each of these three sets is the set of exactly one column, the
column of $M$ has a smaller index than those of $P$ and $Q$, and every
other column with at least three rows avoids $\{a,\dots,g\}$. Columns
with at most two rows may meet these rows. Let $B(A)$ be the number of
columns that are the middle $M$ of an isolated bad path.

\begin{lemma}\label{lem:exchange}
Assume every column has at most three rows, and let $K$ be the number of
gain-$3$ choices of leftmost greedy.
\begin{enumerate}
\item[(a)] Every middle column of an isolated bad path is a gain-$3$ choice
of leftmost greedy, and every other gain-$3$ choice avoids the seven rows
of that path.
\item[(b)] Isolated bad paths with distinct middles have disjoint row sets,
and $3(K+B)\le m$.
\item[(c)] If moreover $H$ holds and $m\ge2$, then
$L\le\lceil(m-K-B)/2\rceil$ and $G-L\ge(B-1)/2$.
\end{enumerate}
\end{lemma}

\begin{proof}
(a) Let $R=\{a,\dots,g\}$. Until greedy first chooses a column meeting
$R$, the middle column has gain $3$, so greedy does not stop before such a
step. At that step the chosen column has gain $3$, the maximum possible, so
it has three rows, meets $R$, and by isolation is one of the three path
columns. All three have gain $3$ and the middle has the smallest index, so
greedy chooses the middle. Afterwards the ends have gain at most $2$, and
every other column with three rows avoids $R$.

(b) Let $M'$ be a second middle, with ends $P',Q'$. By (a), $M'$ avoids
$R$. If $P'$ met $R$, isolation would make $P'$ one of the columns of the
first path, so $P'\subseteq R$; but $P'$ meets $M'$, which avoids $R$. So
$P'$ and likewise $Q'$ avoid $R$. Now replace each of the $B$ middles
among the $K$ gain-$3$ choices by its two ends. The ends of one path are
disjoint, different paths are disjoint, and the remaining gain-$3$ choices
avoid every path by (a), so the resulting $K+B$ triples are pairwise
disjoint and $3(K+B)\le m$.

(c) Cover the $m-3(K+B)$ rows outside these triples in pairs by columns
with exactly those two rows, and a possible last row by any column hitting
it. The middle-only rows $d$ are among them. This gives a cover with
$K+B+\lceil(m-3K-3B)/2\rceil=\lceil(m-K-B)/2\rceil$ columns, and its
indicator vector is feasible for the LP. With Lemma~\ref{lem:phases},
$G-L\ge(m-K)/2-(m-K-B+1)/2=(B-1)/2$.
\end{proof}

Since one column per row gives $L\le m$, part (c) yields: if $H$ holds,
$m\ge2$, $\alpha\ge0$, $B\ge\alpha m$ and $\alpha m\ge2$, then
\begin{equation}\label{eq:gap}
 G-L\ge\frac{\alpha m-1}2\ge\frac{\alpha m}4\ge\frac{\alpha L}4 .
\end{equation}

\section{Counting paths}

\begin{lemma}\label{lem:count}
With parameters \eqref{eq:params} and $k\ge14$, $\mathbb E\,B\ge ck$.
\end{lemma}

\begin{proof}
Call an ordered $7$-tuple $t=(a,\dots,g)$ of distinct rows
\emph{occurring} if $\{a,b,c\},\{c,d,e\},\{e,f,g\}$ are the sets of an
isolated bad path. Each middle column is represented by exactly $8$
occurring tuples: its ends are determined by isolation, and the
representations are given by the choice of which end is $P$ and the orders
of the two private rows of each end. Hence $\sum_t\mathbf 1[t\text{
occurs}]=8B$ on every matrix. By Lemma~\ref{lem:path} in the appendix,
each of the $k(k-1)\cdots(k-6)\ge k^7/128$ tuples occurs with probability
at least $1/(192\cdot10^9k^6)$. By linearity, $\mathbb E\,B\ge
k/(8\cdot128\cdot192\cdot10^9)=ck$.
\end{proof}

\begin{lemma}\label{lem:tails}
On every matrix, $B\le m$. With parameters \eqref{eq:params},
$\Pr[\text{some }|A_j|\ge4]\le1/(10^4k^2)$, and for $k\ge2$,
$\Pr[\text{some pair of rows is the set of no column}]\le160000/k^2$.
\end{lemma}

\begin{proof}
Replace every column with at least four rows by an empty column in the
same position. Every isolated bad path survives, since its three columns
have three rows and an empty column meets nothing; so $B$ does not
decrease, and Lemma~\ref{lem:exchange}(b) for the new matrix gives $B\le
m/3$, in particular $B\le m$, which is all we use. The two probability bounds are union bounds, computed in
Appendix~\ref{app:prob}.
\end{proof}

\begin{proof}[Proof of Theorem~\ref{thm:leftmost}]
Since $B\le m=k$ by Lemma~\ref{lem:tails}, we have
$B\le ck/2+k\,\mathbf 1[B\ge ck/2]$ pointwise. Taking expectations and
using Lemma~\ref{lem:count}, $ck\le\mathbb E\,B\le ck/2+k\Pr[B\ge ck/2]$,
so $\Pr[B\ge ck/2]\ge c/2$. By
Lemma~\ref{lem:tails}, $H$ fails with probability at most
$(10^{-4}+160000)/k^2\le c/4$ for $k\ge10^{20}$. So $H$ and $B\ge ck/2$
hold together with probability at least $c/4$. On this event $A$ is
feasible, and \eqref{eq:gap} with $\alpha=c/2$, valid because $ck/2\ge2$,
gives $G-L\ge cL/8$.
\end{proof}

\section{Uniform tie-breaking}\label{sec:uniform}

For a permutation $\pi$ of the columns, let $A^\pi$ be the matrix with
columns $A^\pi_j=A_{\pi(j)}$. The law of $A$, the event $H$, feasibility
and $L$ are unchanged when $A$ is replaced by $A^\pi$.

\begin{lemma}\label{lem:scan}
Let $(T_\ell)_{\ell\in\Lambda}$ be a family of $3$-element sets indexed
by a finite set $\Lambda$ of labels; different labels may carry the same
set. Consider two processes:
(X) repeatedly choose, uniformly at random, a label whose set is disjoint
from the sets of all labels chosen so far, until none is left; (Y) list
$\Lambda$ in a uniformly random order and accept each label whose set is
disjoint from the sets of the labels already accepted. The sets of chosen
and of accepted labels have the same law.
\end{lemma}

\begin{proof}
Induct on $|\Lambda|$. In (Y), the first accepted label is the first in
the order, uniform on $\Lambda$, as in (X). Given that it is $r$, the
remaining order is uniform on $\Lambda\setminus\{r\}$, labels whose sets
meet $T_r$ (other labels carrying $T_r$ included) are never accepted, and
the order it induces on the set $\Lambda'$ of labels whose sets are
disjoint from $T_r$ is uniform. So both processes continue as the same
processes on $\Lambda'$.
\end{proof}

\begin{proof}[Proof of Theorem~\ref{thm:uniform}]
Assume $H$ and $m\ge2$. While some column has gain $3$, the columns of
maximum gain are the columns with three rows disjoint from all earlier
choices, so the gain-$3$ phase of uniform greedy is process (X) with the
columns with three rows as labels. Leftmost greedy on $A^\pi$ is process (Y) with
the order given by $\pi$, which is uniform when $\pi$ is. By
Lemmas~\ref{lem:phases} and~\ref{lem:scan}, for every predicate $F$ on
integers, the conditional probability given $A$ that $F(\Gu(A))$ holds is
the average of $\mathbf 1[F(G(A^\pi))]$ over all permutations $\pi$. Take
$F(g)$ to be ``$(1+c/8)L(A)\le g$''. Averaging over $A$, using that the law
of $A$, $H$ and $L$ are invariant and substituting $A$ for $A^\pi$,
\[
 \Pr\bigl[H\text{ and }\Gu\ge(1+\tfrac c8)L\bigr]
 =\Pr\bigl[H\text{ and }G\ge(1+\tfrac c8)L\bigr]\ge\frac c4,
\]
where the inequality is the proof of Theorem~\ref{thm:leftmost}. Since $H$
implies feasibility, the theorem follows, with $c/4\ge c/8$.
\end{proof}

\begin{proof}[Proof of Corollary~\ref{cor:conj}]
For $n^{1/10}\ge10$, $\tfrac1{20}n^{1/20}\le\tfrac12n^{1/10}<m(n)\le n$
and $n^{-1/2}\le p(n)\le1/2$, so Assumption~2 holds with $c_0=1/20$, $C_0=1$,
$\delta=1/2$; and $m(n)p(n)\le n^{-3/10}/10\to0$. On the event of
Theorem~\ref{thm:leftmost} the matrix is feasible, $L\ge1$ and
$G/L-1\ge c/8$, whatever values are assigned on infeasible matrices. Along
$n=k^{10}$ this event has probability at least $c/4$, so
$\Pr[|G/L-1|\ge c/8]$ does not tend to $0$. The same argument with
Theorem~\ref{thm:uniform} applies to $\Gu/L$.
\end{proof}

\section{Formal verification}\label{sec:lean}

All of the above is proved in Lean~4 \cite{Lean4} (version 4.19.0) with
Mathlib \cite{Mathlib}, and
the build checks that the main theorems use only the axioms
\texttt{propext}, \texttt{Classical.choice} and \texttt{Quot.sound}. A
matrix is a function giving the set of rows of each column. The
probability of an event is the finite sum, over all matrices, of the
product over columns of $\prod_i(p\text{ or }1-p)$. Greedy is a recursive
function scanning columns by index; $L$ is the real infimum of the
objective over the LP's feasible set, the actual LP value rather than a
bound. Uniform ties are a recursive kernel that averages over the columns
of maximum positive gain at each step. Theorem~\ref{thm:leftmost} is
{\small
\begin{verbatim}
theorem GreedyLPHittingSet.leftmost_explicit : ∀ k ≥ Uniform.threshold,
    (Uniform.c : ℝ) / 4 ≤ matrixProb k (k^10) (parameterP k)
      (fun A => feasible A ∧ (1 + (Uniform.epsilon : ℝ)) * L A ≤ (G A : ℝ))
\end{verbatim}
}
\noindent where the build also checks that \texttt{Uniform.threshold} is
$10^{20}$, \texttt{Uniform.c} is $c$, \texttt{Uniform.epsilon} is $c/8$ and
\texttt{parameterP k} is $1/(10k^4)$. The other statements are:
{\raggedright
\begin{itemize}
\item Theorem~\ref{thm:uniform}:
\texttt{ADG.Uniform.uniform\_probability\_endpoint}, the same statement
with the joint probability of matrix and uniform choices and the bound
$c/8$.
\item Corollary~\ref{cor:conj}, leftmost ties:
\texttt{ADG.full\_frozen\_endpoint}; uniform ties:
\texttt{ADG.Uniform.full\_uniform\_endpoint}. Each states Assumption~2 of
\cite{ADG} as printed, the identities $m(k^{10})=k$ and
$p(k^{10})=1/(10k^4)$, $m(n)p(n)\to0$, and, for every assignment of values
on infeasible matrices, the failure of convergence in probability to one.
\end{itemize}
\par}
Nothing beyond these definitions and the Lean kernel is assumed; the
comparison of the definitions with \cite{ADG} is the one made in
Section~1. Appendix~\ref{app:lean} lists the modules.

To reproduce, run \texttt{lake exe cache get} and \texttt{lake build} in
the \texttt{lean} folder accompanying this paper; with the Mathlib cache
this takes a few minutes.

\begin{appendices}
\section{Probability estimates}\label{app:prob}

Write $\mu(s)=p^{|s|}(1-p)^{m-|s|}$ for the probability that a column has
set $s$; these masses sum to $1$. Under \eqref{eq:params},
$kp\le1/10$.

\begin{lemma}\label{lem:path}
For $k\ge14$ and a fixed tuple $t$ of seven distinct rows,
$\Pr[t\text{ occurs}]\ge1/(192\cdot10^9k^6)$.
\end{lemma}

\begin{proof}
Let $q=p^3(1-p)^{k-3}$ be the mass of each of the three path sets, and $F$
the total mass of the sets $s$ with $|s|\ge3$ and
$s\cap\{a,\dots,g\}\ne\emptyset$; this includes the three path sets. The three path columns occupy three distinct positions with the
middle first, $2\binom n3=n(n-1)(n-2)/3$ choices, and these events are
disjoint because each path set occurs exactly once. So
\[
 \Pr[t\text{ occurs}]=\frac{n(n-1)(n-2)}3\,q^3(1-F)^{n-3}.
\]
A set counted in $F$ contains one of the seven rows and two others, so
$F\le7\binom{k-1}2p^3$ and $nF\le7/1000$. Bernoulli's inequality gives
$q\ge p^3(1-kp)\ge p^3/2$ and $(1-F)^{n-3}\ge1-nF\ge1/2$, and
$n(n-1)(n-2)\ge n^3/4$ since $n\ge4$. Hence the probability is at least
$n^3p^9/192=1/(192\cdot10^9k^6)$.
\end{proof}

\emph{Large columns.} A column with at least four rows contains one of the
$\binom k4$ four-sets, each with probability $p^4$, so
$\Pr[\text{some }|A_j|\ge4]\le n\binom k4p^4\le nk^4p^4=1/(10^4k^2)$.

\emph{Missing pairs.} A column has set exactly $\{a,b\}$ with probability
$r=p^2(1-p)^{k-2}\ge p^2/2$. Its absence from all $n$ columns has
probability $(1-r)^n$. The binomial term $\binom n2r^2(1-r)^{n-2}$ is at
most $1$, so $(1-r)^n\le1/(\binom n2r^2)\le16/(n^2p^4)$, using
$\binom n2\ge n^2/4$. Over at most $k^2$ pairs, the probability is at most
$16k^2/(n^2p^4)=160000/k^2$.

\emph{Threshold.} For $k\ge10^{20}$ we have $ck/2\ge2$, since
$4\cdot196608\cdot10^9\le10^{20}$, and $(10^{-4}+160000)/k^2\le c/4$,
since $4\cdot196608\cdot10^9(10^{-4}+160000)\le10^{40}$.

\section{The Lean modules}\label{app:lean}
{\raggedright

\texttt{Endpoint} and \texttt{SourceEndpoint} contain the definitions and
statements; \texttt{LP} the basic facts on $L$, among them that $L$ is at
most the size of any cover; \texttt{PathDefinitions} the definition of
isolated bad paths and of $B$; \texttt{Probability}, \texttt{PathProbability} and
\texttt{QuantitativeProbability} the finite probability calculations;
\texttt{Greedy} Lemma~\ref{lem:phases}; \texttt{GreedyPaths},
\texttt{DeterministicExchange} and \texttt{DeterministicGap}
Lemma~\ref{lem:exchange} and \eqref{eq:gap}; \texttt{PathCounting} and
\texttt{ExpectationConstruction} Lemma~\ref{lem:count};
\texttt{BadMiddleBound} and \texttt{ExpectationTail}
Lemma~\ref{lem:tails} and the tail bound; \texttt{FullConstruction},
\texttt{SourceTransport} and \texttt{SourceConsequence}
Theorem~\ref{thm:leftmost} and Corollary~\ref{cor:conj}. For uniform ties,
\texttt{UniformEndpoint} defines the kernel and statements;
\texttt{UniformKernel} shows that the kernel gives probability laws;
\texttt{ConflictCoupling} and \texttt{PermutationScan} prove
Lemma~\ref{lem:scan} for an arbitrary directed conflict relation;
\texttt{UniformPhase}, \texttt{OrderedPhase} and \texttt{UniformCount}
identify the gain-$3$ phases; \texttt{UniformInvariance},
\texttt{UniformTransport} and \texttt{UniformGoodLaw} give the permutation
invariance; \texttt{UniformConstruction} proves Theorem~\ref{thm:uniform}.
\texttt{Statement} restates the theorems with their definitions unfolded,
\texttt{Examples} checks small matrices (a single path, where greedy and
the exchange both use three columns, and two paths, where greedy uses six
and the exchange five), and \texttt{Axioms} fails the build if a main
theorem uses another axiom.\par}
\end{appendices}

\begin{thebibliography}{99}
\bibitem{ADG} G. Arpino, D. Dmitriev and N. Grometto,
\emph{Greedy heuristics and linear relaxations for the random hitting set
problem}, in: Approximation, Randomization, and Combinatorial Optimization
(APPROX/RANDOM 2024), LIPIcs 317, Paper 30, 2024.
\href{https://doi.org/10.4230/LIPIcs.APPROX/RANDOM.2024.30}{doi:10.4230/LIPIcs.APPROX/RANDOM.2024.30}.
\bibitem{ADGarxiv} G. Arpino, D. Dmitriev and N. Grometto,
\emph{Greedy heuristics and linear relaxations for the random hitting set
problem}, \href{https://arxiv.org/abs/2305.05565v1}{arXiv:2305.05565v1}, 2023.
\bibitem{BB} D. Bal and P. Bennett, \emph{The matching process and
independent process in random regular graphs and hypergraphs}, Electron.
J. Combin. 30(1) (2023), P1.11.
\href{https://doi.org/10.37236/10698}{doi:10.37236/10698}.
\bibitem{BFPS} J. Blot, W. Fernandez de la Vega, V. Th. Paschos and R. Saad,
\emph{Average case analysis of greedy algorithms for optimisation problems
on set systems}, Theoret. Comput. Sci. 147 (1995), 267--298.
\href{https://doi.org/10.1016/0304-3975(95)00242-O}{doi:10.1016/0304-3975(95)00242-O}.
\bibitem{Chebolu} P. Chebolu, \emph{Topics in random graphs}, PhD thesis,
Carnegie Mellon University, 2008.
\bibitem{Dmitriev} D. Dmitriev, \emph{High-dimensional optimization problems
in average-case and worst-case settings}, PhD thesis, ETH Z\"urich, 2025.
\href{https://doi.org/10.3929/ethz-b-000743377}{doi:10.3929/ethz-b-000743377}.
\bibitem{DFP} M. Dyer, A. Frieze and B. Pittel, \emph{The average
performance of the greedy matching algorithm}, Ann. Appl. Probab. 3 (1993),
526--552. \href{https://doi.org/10.1214/aoap/1177005436}{doi:10.1214/aoap/1177005436}.
\bibitem{FPS} W. Fernandez de la Vega, V. Th. Paschos and R. Saad,
\emph{Average case analysis of a greedy algorithm for the minimum hitting set
problem}, in: LATIN '92, Lecture Notes in Comput. Sci. 583, Springer, 1992,
130--138. \href{https://doi.org/10.1007/BFb0023824}{doi:10.1007/BFb0023824}.
\bibitem{Johnson} D. S. Johnson, \emph{Approximation algorithms for
combinatorial problems}, J. Comput. System Sci. 9 (1974), 256--278.
\href{https://doi.org/10.1016/S0022-0000(74)80044-9}{doi:10.1016/S0022-0000(74)80044-9}.
\bibitem{KS} R. M. Karp and M. Sipser, \emph{Maximum matching in sparse
random graphs}, in: 22nd Annual Symposium on Foundations of Computer
Science (1981), 364--375.
\href{https://doi.org/10.1109/SFCS.1981.21}{doi:10.1109/SFCS.1981.21}.
\bibitem{Lean4} L. de Moura and S. Ullrich, \emph{The Lean~4 theorem
prover and programming language}, in: Automated Deduction -- CADE 28,
Lecture Notes in Comput. Sci. 12699, Springer, 2021, 625--635.
\href{https://doi.org/10.1007/978-3-030-79876-5_37}{doi:10.1007/978-3-030-79876-5\_37}.
\bibitem{Lovasz} L. Lov\'asz, \emph{On the ratio of optimal integral and
fractional covers}, Discrete Math. 13 (1975), 383--390.
\href{https://doi.org/10.1016/0012-365X(75)90058-8}{doi:10.1016/0012-365X(75)90058-8}.
\bibitem{Mathlib} The mathlib Community, \emph{The Lean mathematical
library}, in: Proceedings of the 9th ACM SIGPLAN International Conference
on Certified Programs and Proofs (CPP 2020), 367--381.
\href{https://doi.org/10.1145/3372885.3373824}{doi:10.1145/3372885.3373824}.
\bibitem{MT} M. M\'ezard and M. Tarzia, \emph{Statistical mechanics of the
hitting set problem}, Phys. Rev. E 76 (2007), 041124.
\href{https://doi.org/10.1103/PhysRevE.76.041124}{doi:10.1103/PhysRevE.76.041124};
\href{https://arxiv.org/abs/0707.0189v2}{arXiv:0707.0189v2}.
\bibitem{Slavik} P. Slav\'ik, \emph{A tight analysis of the greedy
algorithm for set cover}, in: Proceedings of the 28th Annual ACM Symposium on
Theory of Computing (STOC 1996), 435--441.
\href{https://doi.org/10.1145/237814.237991}{doi:10.1145/237814.237991}.
Journal version: J. Algorithms 25 (1997), 237--254,
\href{https://doi.org/10.1006/jagm.1997.0887}{doi:10.1006/jagm.1997.0887}.
\bibitem{TZ} O. A. Telelis and V. Zissimopoulos, \emph{Absolute
$o(\log m)$ error in approximating random set covering: an average case
analysis}, Inform. Process. Lett. 94 (2005), 171--177.
\href{https://doi.org/10.1016/j.ipl.2005.02.009}{doi:10.1016/j.ipl.2005.02.009}.
\end{thebibliography}
\end{document}
